\section{Choosing the cheapest action }
\label{sec:rule}
\textbf{Code:} \code{experiments/t1\_expected\_loss.py}.

Suppose a model tells us an email is spam with probability $\post$, and \code{decision.py} has to
pick an action. Delivering it then costs
$\cfn\post$ on average, filing it as spam costs $\cfp(1-\post)$, and inspecting it costs $\ci$
whatever the email is. Figure~\ref{fig:t1} draws these three lines; the cheapest action is the lowest
one. Reading off where the lines cross gives the rule
\[
  \text{deliver if } \post\le\thz=\frac{\ci}{\cfn}=0.3,\qquad
  \text{spam folder if } \post\ge\tho=1-\frac{\ci}{\cfp}=0.97,\qquad
  \text{inspect in between.}
\]
An email goes to the spam folder only if the model is at least $97\%$ sure, because losing a real email
costs ten times as much as delivering a spam. Without inspection the rule is a single cut at
$\tstar=\cfp/(\cfp+\cfn)=10/11\approx0.91$, not the $0.5$ that maximises accuracy. This is the
minimum-expected-loss rule with a reject option~\cite[Section~1.5]{bishop2006}. The flat inspect line
only helps if it lies below the peak where the other two cross, which is at height
$\cstar=\cfp\cfn/(\cfp+\cfn)=10/11$. At $\ci=1$ it does not, and the rule falls back to the single cut
(Figure~\ref{fig:t1}, right).

\begin{figure}[!htbp]
\centering
\includegraphics[width=0.92\linewidth]{t1_expected_loss.pdf}
\caption{Expected cost of delivering, filing as spam and inspecting against $\post$, for
$\cfp=10$, $\cfn=1$. Shading marks where each action is cheapest. Left: $\ci=0.3$. Right: $\ci=1$.}
\label{fig:t1}
\end{figure}

These thresholds are what every later experiment applies. They assume a probability for every
email, so the next step was to build models that give one.
